\documentclass[conference]{IEEEtran}

\usepackage[utf8]{inputenc}
\usepackage[T1]{fontenc}
\usepackage[spanish,es-nodecimaldot]{babel}
\usepackage{amsmath,amssymb}
\usepackage{booktabs}
\usepackage{graphicx}
\usepackage{hyperref}
\usepackage{microtype}
\usepackage{xcolor}
\usepackage{pgfplots}
\pgfplotsset{compat=1.18}

\graphicspath{{figuras/}{build/figuras/}}
\newcommand{\pendiente}[1]{\textcolor{red}{\textbf{PENDIENTE: #1}}}
\renewcommand{\IEEEkeywordsname}{Palabras clave}

\hypersetup{
  pdftitle={Modelo paralelo de segregación residencial basado en Schelling aplicado a Montevideo},
  pdfauthor={Camilo Ferrero y Matías Arce},
  pdfsubject={Simulación secuencial e híbrida MPI/OpenMP del modelo de Schelling},
  pdfkeywords={Schelling, segregación residencial, MPI, OpenMP, HPC},
  hidelinks
}

\title{Modelo paralelo de segregación residencial basado en Schelling aplicado a Montevideo}

\author{
\IEEEauthorblockN{Camilo Ferrero y Matías Arce}
\IEEEauthorblockA{Facultad de Ingeniería, Universidad de la República\\
Montevideo, Uruguay\\
\{camiloferrero13, sonmak44\}@gmail.com}
}

\begin{document}
\maketitle

\begin{abstract}
Este trabajo presenta una simulación basada en el modelo de Schelling para estudiar escenarios de segregación residencial en una representación simplificada de Montevideo. El modelo incorpora composición socioeconómica del vecindario, restricciones de tolerancia, precios de vivienda, capacidad económica y permanencia mínima. Se desarrollaron una implementación secuencial reproducible y una implementación paralela híbrida con MPI y OpenMP, equivalentes bit a bit en las configuraciones verificadas. En una medición local de una iteración sobre una grilla de $1024\times640$, la mediana secuencial fue 3,672~s y la mejor configuración híbrida, con dos procesos y cuatro threads por proceso, alcanzó 1,091~s, un speedup de 3,366. El resultado confirma que la búsqueda de destinos admite paralelismo útil, aunque la eficiencia decrece por trabajo replicado, comunicación y costos seriales. La evaluación entre nodos de FING queda separada de estos resultados locales para evitar extrapolaciones entre plataformas.
\end{abstract}

\begin{IEEEkeywords}
Schelling, modelos basados en agentes, segregación residencial, MPI, OpenMP, computación de alta performance.
\end{IEEEkeywords}

\section{Introducción}

La segregación residencial es un fenómeno emergente en el que numerosas decisiones locales pueden producir distribuciones espaciales globales marcadas. El modelo de Schelling muestra cómo preferencias individuales moderadas pueden generar segregación a escala colectiva \cite{schelling1971}.

Este proyecto adapta el modelo a una representación simplificada de Montevideo e incorpora factores sociales y económicos. Su escala y la necesidad de evaluar múltiples escenarios motivan el uso de computación de alto desempeño.

\subsection{Objetivos y contribuciones}

Las contribuciones previstas son una especificación reproducible del modelo, una implementación secuencial de referencia, una versión híbrida MPI/OpenMP y una evaluación de corrección, desempeño y escalabilidad. \pendiente{ajustar las contribuciones a lo efectivamente implementado.}

\subsection{Motivación computacional}

La evaluación repetida de vecindarios y destinos para cientos de miles de hogares, durante numerosos meses y escenarios, produce un costo suficiente para estudiar paralelismo de memoria distribuida y compartida. \pendiente{cuantificar el costo observado y evitar justificar HPC solamente por el tamaño nominal.}

\subsection{Organización del documento}

Las secciones siguientes describen el modelo y los datos, las implementaciones, la metodología de evaluación, los resultados y las conclusiones.

\section{Descripción del problema y modelo}

\subsection{Grilla, viviendas y hogares}

Montevideo se representa mediante una grilla bidimensional rectangular con celdas residenciales y no residenciales, siguiendo la abstracción espacial del modelo clásico de Schelling. Una celda residencial contiene como máximo una vivienda, que puede estar ocupada o vacía. Los hogares se clasifican en siete subestratos y tres clases socioeconómicas.

El escenario completo inicial utiliza $1024\times640$ celdas. Sus coordenadas expresan proximidad discreta dentro del modelo y no una proyección cartográfica ni ubicaciones reales de Montevideo. Los ocho municipios se representan como zonas contiguas dimensionadas a partir de datos agregados. Esta simplificación permite conservar totales oficiales sin afirmar que la grilla reproduce la forma, las distancias o la trama urbana de la ciudad. La grilla utiliza bordes cerrados: los vecindarios se recortan en los límites y los extremos opuestos no se conectan.

\subsection{Vecindario y ponderación}

Para una celda residencial $c$, el vecindario de radio $r$ se define como
\begin{equation}
N_r(c)=\{c'\in C_R : d(c,c')\leq r\}.
\end{equation}
La influencia disminuye con la distancia mediante un kernel gaussiano,
\begin{equation}
w(c,c')=\exp\left(-\frac{d(c,c')^2}{2\sigma^2}\right).
\end{equation}

El escenario base utiliza pertenencia de Chebyshev con $r=2$, por lo que considera un cuadrado de hasta 24 vecinos, y calcula el peso con distancia euclidiana y $\sigma=1{,}0$. Separar ambas métricas permite un recorrido rectangular eficiente y, a la vez, asignar menor influencia a diagonales y celdas lejanas. Como referencia se evaluará una variante clásica con $r=1$ y ponderación uniforme.

\subsection{Satisfacción y tolerancias}

La proporción ponderada de cada clase se compara con una matriz configurable de preferencias mínimas y tolerancias máximas. La propia unidad familiar no participa del cálculo.

Cuando no existe ningún hogar vecino, la unidad familiar se considera satisfecha por aislamiento. La misma regla se aplica al evaluar una vivienda candidata y estos casos se contabilizan por separado en las métricas. Esta convención evita asignar proporciones artificiales a un denominador vacío.

La configuración base exige para una unidad de clase alta al menos 70\% de vecinos altos, como máximo 30\% medios y 5\% bajos; para una unidad media, como máximo 70\% altos, al menos 50\% medios y como máximo 40\% bajos; y para una unidad baja, como máximo 80\% altos, 90\% medios y al menos 2\% bajos. Todas las condiciones de la fila deben cumplirse simultáneamente.

\subsection{Precios y restricción económica}

El precio propuesto sigue un modelo hedónico simplificado:
\begin{equation}
\log p_v(t)=\alpha_0+\beta_1D_v(t)+\beta_2A_v(t)+\epsilon_v(t).
\end{equation}
La demanda $D_v(t)$ es la proporción ponderada de viviendas residenciales ocupadas del vecindario. El poder adquisitivo $A_v(t)$ es el ingreso medio ponderado de sus hogares, normalizado entre los ingresos de referencia de B$-$ y A+. Ambas variables pertenecen a $[0,1]$; cuando no existe un denominador válido se asigna cero. El ruido $\epsilon_v(t)$ sigue una normal con media cero y desviación $0{,}05$, truncada a $[-0{,}15;0{,}15]$, y se genera de forma determinista por vivienda e iteración mediante un generador sin estado compartido.

El precio y los ingresos se expresan en miles de pesos constantes de 2022. Para aproximar la cuota se usa una anualidad de cuota fija sobre el 80\% del precio, con tasa nominal anual de 6\% y plazo de 240 meses. Estos parámetros son configurables y no pretenden reproducir un producto hipotecario real. Una vivienda es accesible cuando su cuota no supera la fracción $\rho$ del ingreso del hogar.

La configuración base adopta $\beta_1=0{,}40$, $\beta_2=0{,}60$ y $\rho=0{,}30$. El intercepto no se selecciona arbitrariamente: se calibra para que un hogar M, en un entorno con $D=0{,}90$ y el poder adquisitivo normalizado correspondiente a su ingreso, pague una cuota equivalente al 25\% de dicho ingreso. Con los supuestos financieros anteriores se obtiene un precio de referencia de $1448{,}15$ miles de pesos de 2022 y $\alpha_0=6{,}793886203$. Esta cuota objetivo deja un margen respecto del límite de accesibilidad del 30\%.

\subsection{Mudanzas y evolución temporal}

La simulación es síncrona. Los hogares insatisfechos y no bloqueados solicitan destinos accesibles; los conflictos se resuelven de forma determinista antes de aplicar simultáneamente las mudanzas. Una vivienda liberada queda disponible en la iteración siguiente.

Después de mudarse en el mes $t$, un hogar permanece bloqueado durante los 12 meses siguientes y vuelve a ser elegible en $t+13$. Su satisfacción se continúa evaluando durante ese período. El parámetro es configurable y se considera una variante sin permanencia mínima.

\section{Datos, preparación e hipótesis}

La distribución socioeconómica se basa en los subestratos del INSE y la cantidad de viviendas y hogares se aproxima a partir del Censo 2023 \cite{inse2023,censo2023}.

Para Montevideo, la Tabla D3 del INSE 2023 informa las proporciones de hogares B$-$ 5,0\%, B+ 10,4\%, M$-$ 16,0\%, M 21,8\%, M+ 21,2\%, A$-$ 17,1\% y A+ 8,6\%. Estas cifras sustituyen en el escenario base a la distribución preliminar no específica de Montevideo. Los ingresos de referencia continúan siendo los promedios per cápita por subestrato con valor locativo estimados sobre la ECH 2022.

\subsection{Fuentes y transformaciones}

El catálogo de microdatos del INE identifica el estudio como \texttt{URY-INE-CPHV-2023}, con fecha de referencia 31 de mayo de 2023 y unidad geográfica mínima de segmento censal \cite{censo2023ponderado}. La edición de mayo de 2026 incorporó ponderadores para corregir una omisión censal estimada en 10,3\%; en julio de 2026 el INE publicó una revisión posterior del mismo estudio. El catálogo enumera los archivos \texttt{personas\_ext\_05\_2026\_extraccion} y \texttt{viviendas\_ext\_05\_2026\_extraccion}.

El prototipo no descarga, redistribuye ni procesa esos microdatos individuales. Por tanto, no existen archivos censales dentro del repositorio, checksums que registrar ni una licencia que permita presentarlos como parte del proyecto. Se conservan únicamente la identificación y URL del catálogo oficial, cuya consulta está sujeta a los términos del INE: uso científico o estadístico, publicación agregada y prohibición de reidentificación y redistribución sin consentimiento. La implementación usa agregados publicados y parámetros explícitos en archivos de configuración para generar instancias sintéticas reproducibles.

La información municipal de la Intendencia basada en ECH 2023 se utilizará para indicadores socioeconómicos. Sus tablas de cantidades de viviendas corresponden a los censos 2004 y 2011 y no se tratarán como datos de vivienda de 2023.

\subsection{Generación de instancias}

Además del escenario de Montevideo se utilizan instancias sintéticas pequeñas y medianas para verificar reglas e invariantes. Toda generación registra la semilla y los parámetros efectivos.

El generador recorre dos veces la grilla. En la primera pasada determina de forma reproducible cuáles celdas son residenciales y cuáles están ocupadas, obteniendo así la cantidad exacta de hogares para reservar arreglos contiguos. En la segunda construye celdas y hogares, asigna zonas mediante franjas verticales y selecciona subestratos con los pesos configurados. El escenario base usa proporción residencial 0,885 y ocupación 0,90; los conteos efectivos dependen de la semilla y se informan al iniciar cada ejecución.

\subsection{Supuestos y limitaciones de los datos}

No se imputan registros censales ni se realiza una agregación por municipio. Las siete clases socioeconómicas se sortean con las proporciones agregadas para Montevideo del INSE y las zonas del modelo se construyen como franjas verticales de la grilla; estas zonas son particiones computacionales y no municipios reales. Los valores ausentes en las fuentes no se completan: las tolerancias, precios y reglas económicas sin estimación oficial se declaran como supuestos del escenario. Incorporar microdatos ponderados, geometría censal y una correspondencia reproducible entre segmentos y municipios queda fuera del alcance de esta implementación.

La geometría real no se utiliza en el escenario base. Esta convención elimina una dependencia cartográfica, pero impide interpretar posiciones o distancias simuladas como posiciones o distancias reales. La cartografía censal publicada en mayo de 2026, que se encontraba bajo revisión técnica del INE al diseñar el proyecto, queda reservada para una extensión futura.

\section{Diseño e implementación secuencial}

La versión secuencial define la semántica de referencia y proporciona el tiempo base para calcular speedup.

El esqueleto implementado separa la interfaz de línea de comandos, la configuración y el registro en módulos comunes bajo \texttt{src/common}. El ejecutable \texttt{schelling\_seq} carga valores predeterminados, aplica el archivo de escenario, incorpora las sobrescrituras indicadas por argumentos y valida la configuración efectiva antes de iniciar el modelo. Esta organización permite que el ejecutable híbrido reutilice exactamente la misma interpretación de entradas.

\subsection{Estructuras de datos}

La grilla y los hogares se almacenan en arreglos contiguos separados. El identificador global de una celda coincide con su posición en el arreglo y sigue un recorrido por filas; por ello las coordenadas se derivan mediante división y resto sin almacenarlas en cada celda. Cada celda conserva tipo, zona, precio e identificador del hogar ocupante. El valor centinela $-1$ representa una vivienda vacía, evitando mantener un segundo indicador de ocupación que pudiera contradecirlo.

Los hogares conservan identificador, celda actual, subestrato, clase derivada, ingreso, bloqueo y satisfacción. Al crear o validar un estado se comprueba la correspondencia bidireccional entre hogar y celda, que una celda no residencial nunca esté ocupada y que cada hogar tenga exactamente una vivienda.

\subsection{Funciones y organización del código}

El módulo de modelo administra memoria, coordenadas, definición de celdas, ubicación de hogares e invariantes. El generador sintético construye el estado inicial a partir de parámetros y claves aleatorias estables. El módulo de vecindario precalcula una sola vez los desplazamientos de Chebyshev y sus pesos gaussianos. Sobre esas tablas calcula proporciones ponderadas y evalúa las tolerancias cargadas desde el escenario. \pendiente{completar las responsabilidades de simulación, precios, mudanzas y métricas.}

\subsection{Algoritmo}

Cada iteración evalúa satisfacción sobre la ocupación consolidada, actualiza solamente los precios de viviendas vacías y busca un destino para cada hogar insatisfecho no bloqueado. La búsqueda de referencia recorre todas las vacantes, descarta las inaccesibles o insatisfactorias y elige la de menor distancia, con identificador como desempate. Las solicitudes se almacenan sin modificar la grilla. Para cada vivienda se conserva el hogar de mayor ingreso; los empates usan una clave pseudoaleatoria estable y luego el identificador. Finalmente se liberan todos los orígenes ganadores y recién después se ocupan los destinos, por lo que una vivienda liberada no participa hasta el mes siguiente.

Para $H$ hogares, $V$ viviendas vacías y $K$ desplazamientos del vecindario, esta referencia tiene costo dominante $O(HVK)$ en el peor caso. Esta versión prioriza una semántica comprobable; el índice espacial posterior reduce la cantidad de destinos examinados. \pendiente{agregar el pseudocódigo definitivo después de incorporar el índice espacial.}

\subsection{Reproducibilidad y checkpoints}

El generador pseudoaleatorio no mantiene estado mutable. Una mezcla SplitMix64 recibe una clave compuesta por semilla, iteración, identificador de entidad, propósito y número de muestra. Los 53 bits superiores del resultado se convierten a un valor uniforme en el intervalo abierto $(0,1)$. Separar los propósitos de tipo de celda, ocupación, subestrato, ruido y desempate evita que agregar una muestra a una fase desplace las secuencias de las restantes. Existen vectores de prueba fijos para la función de mezcla. \pendiente{documentar la transformación normal y el formato final de checkpoint.}

En ejecuciones normales se genera un checkpoint cada 50 iteraciones y se conservan los dos últimos completos. La escritura usa archivos temporales y un manifiesto final para que un estado parcial no sea reiniciable. En los benchmarks los checkpoints se desactivan para evitar que la E/S altere la comparación de tiempos.

\subsection{Optimizaciones}

\pendiente{documentar optimizaciones únicamente después de medirlas y validarlas contra la referencia.}

\subsection{Convenciones y simplificaciones}

La grilla rectangular es una abstracción y no un mapa de Montevideo. Cada celda residencial representa una vivienda discreta, sin superficie física asociada. Los datos oficiales se conservan como cantidades y distribuciones agregadas por zona. Los índices comienzan en cero, las celdas se ordenan por filas y los bordes son cerrados: todo desplazamiento exterior se descarta.

\subsection{Dificultades de implementación}

\pendiente{describir problemas relevantes, diagnóstico y solución; omitir incidencias triviales que no aporten al análisis.}

\section{Diseño e implementación paralela}

La solución híbrida combina el modelo SPMD de MPI con paralelismo de memoria compartida mediante OpenMP \cite{mpi,openmp}.

\subsection{Descomposición con MPI}

La grilla se divide en subdominios rectangulares. Cada proceso mantiene celdas propias y un halo de ancho $r$. \pendiente{describir la partición final y justificar el balance entre cómputo y comunicación.}

\subsection{Comunicación y migraciones}

Los procesos intercambian halos, información compacta de viviendas vacías, solicitudes y hogares que migran. \pendiente{documentar primitivas MPI, orden de mensajes y volumen comunicado.}

\subsection{Paralelización local con OpenMP}

Los recorridos independientes se distribuyen entre threads. Las escrituras se realizan en posiciones exclusivas o buffers privados y se consolidan determinísticamente. \pendiente{documentar directivas, planificación, afinidad y sincronizaciones finales.}

\subsection{Corrección y determinismo}

La ocupación resultante debe coincidir con la referencia secuencial para una misma entrada y semilla, independientemente de la cantidad de procesos y threads.

\section{Metodología experimental}

\subsection{Plataformas}

El entorno local de desarrollo es WSL2 con Ubuntu 24.04 x86-64 sobre un AMD Ryzen 5 3600 de seis núcleos y doce hilos, 7,7~GiB de memoria y 16~MiB de caché L3. Se utilizó GCC 13.3.0 y Open MPI 4.1.6. El proyecto se construye principalmente con CMake y mantiene un Makefile portable para el entorno académico.

El programa se escribió en C11 y se compiló con advertencias estrictas. Durante el desarrollo también se revisaron el formato, posibles errores de memoria y usos incorrectos de variables. Los comandos principales de compilación y ejecución se conservan en el archivo \texttt{README.md}.

Las mediciones de FING se realizaron el 9 de agosto de 2026 en los equipos
\texttt{pcunix40} y \texttt{pcunix42}. Cada equipo posee un Intel Core i3-4170
a 3,70~GHz, dos núcleos físicos, cuatro hilos lógicos, 15~GiB de memoria y
3~MiB de caché L3. Ambos ejecutaban Fedora Linux con kernel 6.19.14. Se utilizó
GCC 15.2.1 y MPICH 4.2.2, habilitado mediante el módulo
\texttt{mpi/mpich-x86\_64}. Los trabajos se lanzaron directamente por SSH sobre
los dos equipos disponibles, sin un planificador. El equipo \texttt{lulu} se
utilizó únicamente como punto de acceso. MPICH se ejecutó con el proveedor TCP
para utilizar la red Ethernet de las pcunix.

\subsection{Configuraciones y procedimiento}

Para medir escalabilidad fuerte se mantiene fija la instancia y se cambia la cantidad de procesos MPI e hilos OpenMP. Los checkpoints se desactivan y todas las configuraciones deben terminar con el mismo hash.

Para observar el cambio espacial se reconstruye la grilla inicial con la misma semilla y se lee el estado final guardado por la simulación. Cada celda se pinta según su clase social, o como vivienda vacía o celda no residencial. Además se calcula, para cada hogar no aislado, la proporción ponderada de vecinos de su misma clase. El promedio de esas proporciones permite comparar el agrupamiento inicial y final con un único valor. Se usan el mismo radio y los mismos pesos del modelo.

La validación funcional usa tres tamaños: mínimo de $8\times6$, mediano de $256\times160$ y funcional de $1024\times640$. Las primeras mediciones de rendimiento también usaron el escenario funcional. En ellas se ejecutó un calentamiento y cinco repeticiones por configuración, y se informó la mediana.

La instancia principal \texttt{hpc.conf} tiene $4096\times2560$ celdas y 240 iteraciones. Mantiene la distribución social, las preferencias y la regla económica; usa 20\% de viviendas vacías, permanencia cero, financiación de 60\%, plazo de 360 meses y un límite de cuota de 40\% del ingreso. La grilla es una carga sintética y no una estimación de la cantidad de viviendas de Montevideo.

Debido a la duración de la instancia grande, la campaña principal realiza una ejecución de cada configuración. Las mediciones más pequeñas y repetidas se conservan como referencia de variación. Se comparan la versión secuencial y las configuraciones $1\times1$, $1\times2$, $1\times4$, $2\times1$, $2\times2$ y $2\times4$. El primer número indica procesos MPI y el segundo, hilos OpenMP por proceso.

\subsection{Métricas}

Para $T_1$ secuencial y $T_p$ paralelo, se calculan
\begin{equation}
S_p=\frac{T_1}{T_p}, \qquad E_p=\frac{S_p}{p}.
\end{equation}
También se miden tiempos máximos por fase, desbalance $L=\max_i W_i/\overline{W}$, volumen lógico de comunicación y métricas funcionales del modelo. El trabajo $W_i$ es la cantidad de hogares elegibles cuya búsqueda de destino pertenece al proceso $i$.

\subsection{Reproducibilidad}

El script \texttt{scripts/ejecutar\_experimentos.py} crea un directorio por ejecución, inicia las versiones secuencial e híbrida, compara los hashes y produce \texttt{resumen.csv}. Fija \texttt{OMP\_PLACES=cores}, \texttt{OMP\_PROC\_BIND=close} y la cantidad de hilos de cada configuración. Los archivos de la campaña principal se guardan en \texttt{informe/datos/hpc\_final} junto con la configuración necesaria para repetirla.

El proyecto se compila con CMake y cada versión recibe un archivo mediante \texttt{--config}. La ejecución secuencial usa \texttt{schelling\_seq}; la híbrida se inicia con \texttt{mpirun}, indicando procesos e hilos. La salida incluye la configuración efectiva, métricas y tiempos por iteración, datos de comunicación y el estado final. Con \texttt{hpc.conf} y semilla 42, todas las configuraciones finalizaron con el hash \texttt{6b39928d840ed346}.

\section{Validación funcional}

La validación combina pruebas unitarias, escenarios manuales, invariantes de conservación, análisis de memoria y comparación de hashes por iteración.

\subsection{Casos de prueba}

\pendiente{resumir la batería implementada y su cobertura.}

\subsection{Equivalencia entre versiones}

\pendiente{presentar configuraciones comparadas, tolerancias numéricas y resultados.}

\subsection{Checkpoint y reinicio}

\pendiente{mostrar que una ejecución reiniciada coincide con una ejecución ininterrumpida.}

\section{Resultados}

\subsection{Desempeño local}

La Tabla~\ref{tab:resultados-locales} resume cinco repeticiones medidas después de un calentamiento. Todas produjeron el hash \texttt{4eac656015050f40}.

\begin{table}[ht]
\centering
\caption{Escalabilidad fuerte local para una iteración base}
\label{tab:resultados-locales}
\begin{tabular}{rrrrr}
\toprule
MPI & OMP & $T$ (s) & $S_p$ & $E_p$ \\
\midrule
-- & -- & 3,672 & 1,000 & 1,000 \\
1 & 1 & 3,706 & 0,991 & 0,991 \\
1 & 2 & 1,870 & 1,964 & 0,982 \\
1 & 4 & 1,890 & 1,943 & 0,486 \\
2 & 2 & 1,876 & 1,957 & 0,489 \\
2 & 4 & 1,091 & 3,366 & 0,421 \\
\bottomrule
\end{tabular}
\end{table}

\begin{figure}[ht]
\centering
\begin{tikzpicture}
\begin{axis}[
    width=\columnwidth,
    height=0.62\columnwidth,
    xlabel={Recursos de cómputo},
    ylabel={Speedup},
    xmin=1,xmax=8,
    ymin=0,ymax=4,
    xtick={1,2,4,8},
    grid=major,
    legend style={font=\scriptsize,at={(0.03,0.97)},anchor=north west}
]
\addplot[dashed,gray] coordinates {(1,1) (2,2) (4,4)};
\addlegendentry{Ideal}
\addplot[mark=*,blue] coordinates {(1,0.991) (2,1.964) (4,1.943) (4,1.957) (8,3.366)};
\addlegendentry{Medido}
\end{axis}
\end{tikzpicture}
\caption{Speedup local respecto a la versión secuencial. Los dos puntos de cuatro recursos corresponden a $1\times4$ y $2\times2$.}
\label{fig:speedup-local}
\end{figure}

\subsubsection{Versión secuencial}

La mediana secuencial fue 3,672~s. La muestra máxima alcanzó 5,979~s, mientras las otras cuatro quedaron entre 3,631 y 3,740~s; por ello se usa la mediana como estadístico principal.

\subsubsection{Versión híbrida}

El costo del ejecutable híbrido con un recurso fue 0,9~\% mayor que el secuencial. Dos threads alcanzaron un speedup de 1,964, pero cuatro threads dentro de un solo proceso no mejoraron ese valor. La configuración $2\times4$ obtuvo el menor tiempo, 1,091~s, y un speedup de 3,366. En ella la búsqueda consumió 0,906~s, la preparación 0,057~s, los índices 0,086~s, la comunicación 0,041~s y la consolidación 0,019~s.

Con dos procesos el desbalance medio de trabajo fue 1,002, consistente con la distribución sintética casi uniforme. Se contabilizaron 4373 solicitudes remotas y 23,2~MB lógicos comunicados durante la iteración.

Una pasada adicional de los 120 meses mostró un comportamiento diferente. La versión secuencial demoró entre 32 y 36~s. La configuración $2\times4$ demoró entre 19,8 y 26,0~s; su mejor speedup observado fue 1,63. Esta comparación tiene una sola muestra por configuración y se usa para estimar la duración, no como resultado final. La menor mejora se explica porque, una vez que dejan de existir destinos válidos, las fases comunes y la comunicación tienen mayor peso que la búsqueda paralela.

\subsection{Desempeño en FING}

La Tabla~\ref{tab:resultados-fing} muestra las medianas de cinco repeticiones
del escenario completo de 120 meses. Todas las muestras produjeron el hash
\texttt{9de3a86b1efdd900}, igual al de la versión secuencial.

\begin{table}[ht]
\centering
\caption{Escalabilidad fuerte en FING para 120 meses}
\label{tab:resultados-fing}
\begin{tabular}{rrrrrr}
\toprule
MPI & OMP & $T$ (s) & $S_p$ & $E_p$ & Comunicación \\
\midrule
-- & -- & 28,584 & 1,000 & 1,000 & -- \\
1 & 1 & 32,141 & 0,889 & 0,889 & 0 \\
1 & 2 & 22,636 & 1,263 & 0,631 & 0 \\
1 & 4 & 20,623 & 1,386 & 0,347 & 0 \\
2 & 1 & 35,551 & 0,804 & 0,402 & 2,79 GB \\
2 & 2 & 34,758 & 0,822 & 0,206 & 2,79 GB \\
2 & 4 & 32,884 & 0,869 & 0,109 & 2,79 GB \\
\bottomrule
\end{tabular}
\end{table}

\begin{figure}[ht]
\centering
\begin{tikzpicture}
\begin{axis}[
    width=\columnwidth,
    height=0.62\columnwidth,
    xlabel={Recursos de cómputo},
    ylabel={Speedup},
    xmin=1,xmax=8,
    ymin=0,ymax=2,
    xtick={1,2,4,8},
    grid=major,
    legend style={font=\scriptsize,at={(0.97,0.97)},anchor=north east}
]
\addplot[dashed,gray] coordinates {(1,1) (2,2)};
\addlegendentry{Ideal}
\addplot[mark=*,blue] coordinates {(1,0.889) (2,1.263) (4,1.386)};
\addlegendentry{Un nodo}
\addplot[mark=square*,red] coordinates {(2,0.804) (4,0.822) (8,0.869)};
\addlegendentry{Dos nodos}
\end{axis}
\end{tikzpicture}
\caption{Speedup en FING respecto a la versión secuencial.}
\label{fig:speedup-fing}
\end{figure}

El menor tiempo fue 20,623~s con un proceso y cuatro hilos, un speedup de
1,386. Dos hilos, que corresponden a los dos núcleos físicos del equipo,
alcanzaron 22,636~s. Los cuatro hilos aprovecharon además los hilos lógicos del
procesador, pero la mejora adicional fue menor.

Las configuraciones de dos procesos fueron más lentas que la secuencial. En
los 120 meses intercambiaron 2,79~GB lógicos y atendieron 195391 solicitudes
remotas. El desbalance medio fue solamente 1,005, por lo que la diferencia no
se explica por una mala división del trabajo. El costo proviene principalmente
del intercambio y de mantener datos globales en ambos procesos.

\subsection{Comportamiento del modelo}

En la primera iteración base, 291147 de 522351 hogares estaban satisfechos (55,7~\%). Se generaron 222341 solicitudes, de las cuales 28258 fueron aceptadas y 194083 rechazadas por conflicto; otros 8863 hogares insatisfechos no hallaron destino viable. En el mes 14 se alcanzaron 333477 hogares satisfechos (63,8~\%) y no se produjo ninguna nueva solicitud. Desde ese punto hasta el mes 120, 188874 hogares permanecieron insatisfechos porque ninguna vivienda vacía cumplía a la vez sus preferencias y su capacidad de pago. Las viviendas se liberaron correctamente; el estado quedó trabado por las reglas del escenario sintético.

\subsection{Sensibilidad}

En el escenario mediano se compararon diez iteraciones base, sin ruido de precios y sin permanencia mínima. Los hogares satisfechos finales fueron 20955, 20941 y 20966, respectivamente, sobre 32691 hogares: una variación menor a 0,1 puntos porcentuales. Sin permanencia no quedaron hogares bloqueados, como exige la definición. Las solicitudes y hogares sin destino variaron ampliamente entre escenarios en la décima iteración, indicando que una observación mensual aislada no caracteriza la dinámica; el análisis final deberá usar series completas y varias semillas.

\section{Discusión}

\subsection{Interpretación de resultados}

La búsqueda de destinos domina el tiempo incluso después del índice espacial. El salto de uno a dos threads casi reduce a la mitad esa fase, mientras que cuatro threads en un único proceso no agregan mejora, posiblemente por ancho de banda, planificación del entorno virtualizado o por el trabajo serial de construcción de índices. La configuración $2\times4$ vuelve a reducir la búsqueda al repartir hogares primero por franjas MPI y luego entre threads.

Al agregar recursos, una parte del trabajo sigue repitiéndose en todos los procesos y también aumenta el intercambio de datos. Por eso el tiempo no disminuye en la misma proporción. Los 23,2~MB comunicados con dos procesos son pequeños frente al costo de búsqueda en esta máquina, pero crecerán con la grilla y la cantidad de procesos. Las mediciones en FING permitirán saber cuánto mejora el programa al usar más de una máquina.

El desbalance de 1,002 no justifica por ahora reemplazar las franjas uniformes por cortes dinámicos según población. Como extensión se priorizó el paralelismo entre escenarios: un coordinador ejecuta configuraciones independientes concurrentemente, limita la cantidad de trabajos y combina sus resúmenes. Esta capa resulta útil para sensibilidad y múltiples semillas sin introducir comunicación entre simulaciones.

\subsection{Factores que limitan los resultados}

Los resultados dependen de la grilla simplificada, los datos sociales agrupados, los valores elegidos para las preferencias, la variación de los tiempos y las diferencias entre equipos.

En particular, la grilla rectangular conserva relaciones abstractas de vecindad, pero no la forma urbana, la red vial, las distancias ni la distribución geográfica exacta de viviendas de Montevideo. Esta decisión mejora la claridad algorítmica y facilita la partición paralela, a costa de limitar la interpretación territorial de los resultados.

\subsection{Limitaciones}

La implementación MPI conserva catálogos globales de hogares y vacantes para priorizar corrección y determinismo; por tanto, su consumo de memoria por proceso no escala con el subdominio local. La medición local usa una única iteración y comparte recursos con el sistema anfitrión. Los resultados no deben extrapolarse a FING ni interpretarse como evidencia empírica sobre segregación urbana real.

\section{Conclusiones y trabajo futuro}

Se completó un simulador del modelo extendido de Schelling que permite repetir una ejecución con los mismos resultados. La versión secuencial sirve como referencia y la versión paralela reparte franjas de la grilla con MPI y usa OpenMP para calcular satisfacción, precios y destinos. Ambas versiones produjeron el mismo estado con distintas cantidades de procesos e hilos. También se comprobó que no se pierdan hogares ni viviendas vacías y que una ejecución guardada pueda continuar correctamente.

Las mediciones mostraron que OpenMP acelera la búsqueda cuando existe suficiente trabajo. En la grilla principal, la versión secuencial demoró 3~h 49~min y la mejor configuración, con un proceso y cuatro hilos, demoró 2~h 6~min: una mejora de 1,817. Dos hilos lograron 1,724 con mayor eficiencia. Agregar una segunda máquina no mejoró el mejor tiempo porque cada proceso conserva datos globales y las configuraciones de dos nodos intercambiaron aproximadamente 83,2~GB.

La grilla principal usa permanencia cero, financiación de 60\% y 20\% de viviendas vacías. La satisfacción pasó de 54,6\% a 78,4\% y se aceptaron casi 2,5 millones de mudanzas. La proporción media de vecinos de la misma clase aumentó de 43,7\% a 59,3\%, y la grilla final mostró agrupamientos más claros. Estos resultados coinciden con el comportamiento esperado del modelo de Schelling, dentro de las limitaciones de una grilla sintética.

\subsection{Trabajo futuro}

La versión MPI podría repartir las viviendas vacías entre los procesos en lugar de mantener una copia completa en cada uno. Esto reduciría la memoria y los datos intercambiados, aunque obligaría a coordinar la búsqueda y los conflictos entre procesos. También sería útil repetir la grilla principal con varias semillas y varias ejecuciones por configuración. Finalmente, usar un mapa y datos censales por zona permitiría representar Montevideo con mayor detalle.


\bibliographystyle{IEEEtran}
\bibliography{referencias}

\end{document}
